Suppose that for a given environment, we have specified a probability measure $P$.
In particular, we are interested in the distribution it induces on a random variable $X$, denoted $P(X)$.

Now, imagine we ask an expert---a person familiar with the environment---to evaluate our choice of $P$.
The expert may hold different beliefs and instead favor another probability measure $Q$.
From the expert’s perspective, evaluating $P$ involves assessing the potential loss incurred by using
the (in their view) incorrect distribution $P(X)$ when the true distribution is believed to be $Q(X)$.

Let us assume that a real number can be used to measure this loss, and denote it as $\kl[][][][ ]$.
It seems reasonable, as $P$ gets closer to $Q$ this measure decreases, and increases when $P$ diverges from $Q$.
Therefore, $\measure{D}{}{}$ effectively is a measure of the divergence of $P(X)$ from $Q(X)$, when $Q$ is
the believed true probability measure.

Interestingly, if we make certain assumptions about the properties of this divergence measure,
it turns out that---up to a constant factor---it must be equal to the Kullback–Leibler (KL) divergence.

Alternatively, by adopting a more minimal set of properties, we can prove that if we accept Shannon’s entropy as a
measure of the uncertainty of a distribution, we should also consider KL divergence as the appropriate method for
measuring divergence between distributions.


\begin{theorem}
    Let $X$ be a discrete random variable.
    Suppose that $\kl[][][][ ]$ is some divergence measure that quantifies the divergence of $P(X)$ from $Q(X)$.
    Consider the following properties:

    \begin{description}
        \item[Consistency with entropy] Let $\varOmega$ be a continuous random variable, and $Q$ be a probability
        measure.
        If $U(\varOmega)$ is the uniform distribution with the same support as $Q(\varOmega)$, then we have
        \[
            \kl[Q][U][\varOmega][ ] = \entropy{U}{\varOmega} - \entropy{Q}{\varOmega},
        \]
        where $\measure{H}{}{}$ denotes Shannon's entropy.
        \item[Additivity] For any two discrete random variables $X, Y$ and probability measures $P,Q$
        \[
            \kl[][][X, Y][ ] = \kl[][][][ ] + \sum_x \kl[][][Y|x][ ]Q(x)\cdot
        \]
    \end{description}
    If $\measure{D}{}{\,\cdot\,}$ satisfies these properties, then it is the Kullback-Leibler divergence.
\end{theorem}

\begin{proof}
    \newcommand{\D}[1][]{\kl[][][#1][ ]}
    \newcommand{\T}{\varOmega}
    \renewcommand{\t}{\omega}
    \renewcommand{\tIx}{\t \mid x}
    \newcommand{\X}{X}
    \newcommand{\Y}{Y}

    \newcommand{\TIx}{\T|x}
    \newcommand{\TX}{\T, X}
    \newcommand{\tx}{\t, x}
    \newcommand{\tent}[1]{\int q(#1) \ln q(#1) d\t}
    \newcommand{\Qx}{Q(x)}
    \newcommand{\sumx}{\sum_{x}}

    Let $P(X)$ and $Q(X)$ be two arbitrary distributions for a discrete random variable X.
    We define a new continuous random variable $\T$ and let
    \begin{equation*}
        p(\T = \t \mid \rvEq) =
        \begin{cases}
            \frac{1}{P(x)} & 0 \leq \t \leq P(\rvEq), \\
            0                       & \text{otherwise}\cdot
        \end{cases}
    \end{equation*}
    Notice that both $P(\T|X)$ and $P(\TX)$ are uniform distributions, and
    $\entropy{P}{\TIx} = \ln P(x)$, $\entropy{P}{\TX} = 0$.

    We define $Q(\TIx)$ to be an arbitrary distribution that has the same support as $P(\TIx)$.
    This ensures that $Q(\TX)$ and $P(\TX)$ have the same support as well.
    Thus, we have
    \begin{align*}
        \D[\TX] & = \entropy{P}{\TX} - \entropy{Q}{\TX}    \\
        & = \sumx \tent{\tx}                       \\
        & = \sumx \Qx \int q(\tIx) \ln q(\tx) d\t,
    \end{align*}
    and
    \begin{align*}
        \D[\TIx] & = \entropy{P}{\TIx} - \entropy{Q}{\TIx} \\
        & = \ln P(x) + \tent{\tIx}\cdot
    \end{align*}
    Now, we use the additivity property to obtain a formula for our divergence measure:
    \begin{align*}
        \D & = \D[\TX] - \sumx \D[\TIx]\Qx                                                              \\
        & = \sumx \Qx  \left( \int q(\tIx) \ln \frac{q(\tx)}{q(\tIx)}d\t - \ln P(x) \right) \\
        & = \sumx \Qx  \left( \ln \Qx \int q(\tIx) d\t - \ln P(x) \right)                   \\
        & = \sumx \Qx \ln \frac{\Qx}{P(x)}\cdot \qedhere
    \end{align*}
\end{proof}

In this theorem, the first property simply states that the entropy of a distribution should be related to the
divergence of the uniform distribution from it.
The more the divergence is, the less the entropy must be.

The second property tells us that the divergence is an additive quantity.
We know that the divergence of $P(XY)$ from $Q(XY)$ is a measure of the generic loss of using $P(XY)$
when we believe $Q(XY)$ is the true distribution.
To calculate this loss, we can assume that $X$ happens before $Y$.
In that way, this loss would include the loss of using $P(X)$, and then, if we observe
$\assign{X}$, which can happen with the probability $Q(x)$, it would also include the loss of
using $P(Y \mid \assign{x})$.
This property indicates that for combining these individual terms we should use addition.

\newcommandx{\klRatio}[2][1=Q, 2=R]{\sum_x #1(x) \ln \frac{#2(x)}{P(x)}}

\begin{lemma}
    Let $X$ be a random variable, and let $P$, $Q$, and $R$ be probability distributions over $X$.
    The Kullback–Leibler divergence satisfies the following inequality:
    \begin{equation*}
        \kl \geq \klRatio \cdot
    \end{equation*}
\end{lemma}

\begin{proof}
    $R$ is a probability distribution, and the Kullback-Leibler divergence is non-negative:
    \begin{equation*}
        \kl[][R] \geq 0 \cdot
    \end{equation*}
    By adding $\klRatio$ to both sides of the above inequality, the desired inequality follows.
\end{proof}

\begin{theorem}
    Let $Q_1$, $Q_2$ and $P$ be probability distributions.
    If for some $0 \leq \alpha \leq 1$ we have $Q = \convex{Q_1}{Q_2}$, then:
    \begin{equation*}
        \kl \leq \convex{\kl[Q_1]}{\kl[Q_2]} \cdot
    \end{equation*}
\end{theorem}

